\section{PPO 方法与实验}
\label{sec:ppo}
本节按与 DQN 相同的层次介绍 PPO，重点解释同策略更新、实现变体、开发选模与配对对照。除已有 GAE 消融外，补充 KL 学习率反馈的三次配对续训、历轮改动的收益区间及控制过程诊断。历轮开局集合不同，分数分别按各轮协议解释；DQN 与 GAE 阶段冻结候选的同开局 50 局结果集中见第\ref{sec:comparison}节。

\subsection{算法原理}
\subsubsection{PPO 的策略与价值学习}
与 DQN 估计各动作的 Q 值不同，近端策略优化（PPO）直接学习动作分布 $\pi_\theta(a\mid s)$，并用状态价值网络 $V_\phi(s)$ 辅助评价动作。策略网络输出四个动作的 logits，经 softmax 得到概率；训练时按概率采样，确定性评估时选择概率最大的动作。价值网络输出一个标量，估计当前状态下的未来折扣回报。

PPO 属于同策略方法。它先用旧策略收集一段 rollout，记录状态、动作、奖励、价值预测与旧动作对数概率，再在这批数据上进行若干轮小批量更新。该批更新后继续采集新轨迹，不把任意旧回放直接代入 PPO 概率比目标。它不使用 DQN 的目标 Q 网络或 $\varepsilon$-greedy 探索；探索由动作分布和熵奖励共同调节。

\subsubsection{GAE 优势估计与信用分配}
定义价值时间差分误差与广义优势估计为
\begin{align}
 \delta_t^V &= r_t+\gamma(1-d_t)V_{\phi_{\mathrm{old}}}(s_{t+1})-V_{\phi_{\mathrm{old}}}(s_t),\\
 \widehat A_t &= \sum_{l=0}^{K-1}(\gamma\lambda)^l\delta_{t+l}^V,
 \qquad \widehat R_t=\widehat A_t+V_{\phi_{\mathrm{old}}}(s_t).
\end{align}
其中 $d_t$ 只在真正终止时为 1；$K$ 为从 $t$ 到本局终止、时间截断或当前 rollout 边界前的有效转移数，不跨越环境重置累加优势。真正终止不自举，时间截断则用实际终止观测的价值自举。$\lambda$ 调节偏差和方差，$\widehat R_t$ 是固定的价值回归目标。训练小批量内的优势标准化不改变评估的原始得分。

GAE 对照阶段固定 $\gamma=0.995$，仅比较 $\lambda=0.98$ 与 0.99。权重衰减尺度 $1/(1-\gamma\lambda)$ 约从 40 增至 67 个决策步，有望将较远的击砖奖励联系到更早动作，但也可能增加优势方差。这是权重尺度，不是超过该步数就完全忽略回报；价值自举仍参与长期估计。

\subsubsection{概率比裁剪与训练目标}
新旧策略对已采样动作的概率比为
\begin{equation}
 \rho_t(\theta)=\frac{\pi_\theta(a_t\mid s_t)}{\pi_{\theta_{\mathrm{old}}}(a_t\mid s_t)}.
\end{equation}
PPO 的裁剪代理目标及本项目最小化的总损失为
\begin{align}
 L^{\mathrm{clip}}(\theta)&=\mathbb E_t\!\left[\min\!\left(\rho_t\widehat A_t,
 \operatorname{clip}(\rho_t,1-\epsilon,1+\epsilon)\widehat A_t\right)\right],\\
 \mathcal L(\theta,\phi)&=-L^{\mathrm{clip}}(\theta)
 +c_V\mathbb E_t\!\left[\ell_V\bigl(\widehat R_t-V_\phi(s_t)\bigr)\right]
 -c_H\mathbb E_t\!\left[\mathcal H(\pi_\theta(\cdot\mid s_t))\right].
\end{align}
其中 $\epsilon$ 为概率比裁剪范围，$c_V$ 和 $c_H$ 控制价值回归和熵奖励。裁剪限制继续放大某些概率变化的收益，但不保证每次更新都提高游戏得分；项目另设近似 KL 保护。独立梯度裁剪和稳健价值回归属于 PPO 的实现变体。

\subsection{模型实现与训练管理}
\subsubsection{PPO 策略与价值双网络}
初始标准 PPO 共享卷积特征，后续改进版本将策略与价值特征提取器分开。每个分支均采用 NatureCNN：三层卷积依次使用 32 个 $8\times8$、64 个 $4\times4$ 和 64 个 $3\times3$ 卷积核，步长为 4、2、1，展平后连接 512 维特征，与表\ref{tab:network}的特征提取部分一致；策略头为 $512\rightarrow4$ logits，价值头为 $512\rightarrow1$ 标量。拆分时复制原权重，使迁移前后的动作概率和价值预测相同；参数组改变后重置 Adam，匹配对照做相同重置。推理只使用策略分支。

Balanced PPO 保留标准 rollout、GAE、概率比裁剪和优势标准化，对策略与价值参数分别按范数 0.5 和 1.0 裁剪。价值损失为 $2\times\mathrm{Huber}(\delta=10)$：小误差区与 MSE 相同，大误差区减弱极端残差影响。网络分开不自动意味着梯度预算独立；裁剪缩放也不能直接解释为 Adam 参数步长的等比例变化。

具体地，令价值残差 $e=\widehat R-V_\phi(s)$，后续 GAE 与 KL 对照均采用
\begin{equation}
 \ell_V(e)=\begin{cases}
 e^2, & |e|\leq 10,\\
 20|e|-100, & |e|>10.
 \end{cases}
\end{equation}
因此单样本损失对残差的导数绝对值最多为 20；这与参数梯度范数裁剪作用于不同层面。GAE 与 KL 对照未启用价值预测裁剪，避免将其与策略概率比裁剪混淆。

\reportfigure{images/ppo_actor_critic.png}{PPO 策略与价值网络}{相同图像输入分别送入独立 CNN，输出动作分布与状态价值。}{PPO 改进版本的策略/价值结构。训练时采样动作并拟合 GAE 目标，确定性评估只选择策略概率最大的动作。}{fig:ppo-network}

\subsubsection{训练循环与模型更新}
PPO 的更新循环按完整 rollout 组织，后续对照的共同流程如下：
\begin{enumerate}
 \item 用当前策略在 16 个环境中各采集 256 个决策步，保存旧动作概率和旧价值预测，共得到 4096 条转移。
 \item 在局末或 rollout 边界正确自举，计算 GAE 和固定 returns；后续 epoch 不重算旧策略概率或优势目标。
 \item 打乱这批数据，以 256 条转移为一个小批量，最多优化 4 个 epoch；小批量内标准化优势，分别计算策略、价值和熵项。
 \item 分别裁剪策略与价值梯度，按退火后的学习率更新 Adam；若小批量近似 KL 超过目标的 1.5 倍，则提前结束本轮剩余更新。
 \item 在整批原 rollout 上记录更新后的 KL 和裁剪比例，随后采集新轨迹；按开发协议评估并保留候选模型。
\end{enumerate}
一批轨迹样本内的重复优化属于同一轮 PPO 更新，不是跨轮长期回放。训练预算约为 50 万个决策步，按完整 rollout 批次完成采样，因此实际步数会略高于计划预算。

\subsubsection{续训与变体状态保存}
PPO 续训恢复策略、价值与优化器状态，重新采集 rollout，不恢复历史回放。阶段衰减按新增步数计时，各变体均保存相应元数据。DualRate 的额外价值拟合保持原优势、returns 和旧动作概率，不更新策略权重及其 Adam 状态。KL 反馈变体另外保存平滑 KL、学习率倍率和控制计数，只有显式开始新的优化器阶段时才重置控制状态。

PPO 训练入口为 \texttt{train\_ppo.py}，独立裁剪、双学习率与 KL 反馈封装为更新变体。逐局开发种子与候选选择使用共享模块，以保证不同变体的评估口径一致。

\subsection{训练设置与开发选模}
\subsubsection{PPO 训练阶段与单因素对照}
标准 PPO 首轮从随机初始化训练约 50 万步；之后依次研究独立策略/价值 CNN 与退火、大 rollout、长回报、辅助奖励、独立梯度裁剪和稳健价值损失。辅助奖励曾设置接球奖励与失球惩罚，改变训练目标，未作为最终比较方案。未获得收益的实验同样保留。GAE 同起点对照之后，又开展 KL 反馈的多次配对续训；不同测试开局上的历轮分数不拼成提升曲线。

GAE 阶段用新验证集选出的同一个 Balanced 检查点作为起点，采用相同的随机性设置，两组均重置 Adam。主要参数见表\ref{tab:ppo-hyper}，仅 GAE $\lambda$ 不同，均按相同预算完成训练。多个并行环境由同进程向量环境采样，不代表多次独立训练。

\reportfigure{images/ppo_training_curves.png}{PPO 训练和开发诊断}{同一模型起点的两组 GAE 设置，展示真实开发得分与更新日志。}{PPO GAE 对照的实际日志。开发均分与中位数使用相同 32 个逐局种子；训练均分含随机动作，整轮 KL 在固定 rollout 上更新后测量，不能与开发回报混为一谈。}{fig:ppo-curves}

\begin{table}[!htbp]
 \centering\small
 \caption{PPO 的 GAE 阶段共同训练设置}\label{tab:ppo-hyper}
 \begin{tabularx}{.95\linewidth}{Y Y}
 \toprule
 参数 & 设置\\
 \midrule
 单组新增预算 & 约 50 万个决策步\\
 环境数 $\times$ 每环境 rollout & $16\times256=4096$\\
 批大小 / 优化 epoch & 256 / 4\\
 学习率 & $2\times10^{-5}\rightarrow5\times10^{-6}$\\
 折扣因子 / GAE lambda & 0.995 / 0.98 或 0.99\\
 clip range / KL 目标 & $0.06\rightarrow0.03$ / 0.0075\\
 熵系数 / 价值损失权重 & $0.0003\rightarrow0.0001$ / 0.25\\
 策略 / 价值梯度范数上限 & 0.5 / 1.0\\
 价值损失 / Adam & $2\times\mathrm{Huber}(\delta=10)$ / 两组均重置\\
 随机性设置 / 开发评估 & 两组保持一致 / 定期评估 32 局\\
 \bottomrule
 \end{tabularx}
\end{table}

\Needspace{8\baselineskip}
\subsubsection{开发数据与候选选择}
起点选择、训练开发评估与最终测试使用相互分离的开局集合；开发环境按局重置随机性设置，不依赖前局的随机数消耗；最终测试对各候选统一运行 50 局。最终优先选开发均分最佳，只有中位数最佳模型的均分达到其 95\% 且中位数更高时才替换。该阈值是预先约定的取舍，不是显著性检验。GAE 阶段两组均执行完整预算，稳定候选均未满足最终均分约束，实际选中均分最佳模型。

\subsubsection{KL 反馈与三次配对续训设置}
KL 阶段从 GAE 开发流程选中的 $\lambda=0.98$ 模型冻结副本开始，保留原始奖励和上述独立网络、Huber 损失及梯度裁剪。采用三种不同的训练随机性设置，每种设置各运行反馈组和固定计划对照，组成三次配对续训。每次新增约 50 万个决策步，实际均为 503808 步；两组均重置 Adam，执行完整预算。每个检查点评估相同的 32 局开发开局，最终所选模型各做 50 局独立测试。此阶段另用新的开发和测试开局集合，不与第\ref{sec:comparison}节的原统一测试混合。

基础学习率仍由 $2\times10^{-5}$ 退火到 $5\times10^{-6}$，仅反馈组的策略分支乘以可调倍率：
\begin{equation}
 \eta_{\pi,k}=m_k\eta_{\mathrm{base},k},\qquad
 \eta_{V,k}=\eta_{\mathrm{base},k},\qquad
 m_k\in[0.5,2].
\end{equation}
控制器观测整轮更新后的采样 KL，以 0.2 为新观测权重进行指数平滑。反馈目标为 $5\times10^{-4}$，平滑值低于目标的 $1/1.5$ 时将下一轮倍率乘以 1.1，高于目标的 1.5 倍时除以 1.1；瞬时 KL 超过上界时优先下调，最后限制倍率范围。对照使用同一实现关闭反馈，倍率恒为 1，价值学习率及其余设置一致。该反馈目标与小批量 KL 保护参数 0.0075 作用不同：前者调整下一轮策略学习率，后者限制当前更新。

\subsection{优化结果与负结果分析}
标准 PPO 在首轮 20 局测试中均分为 35.00；稳定续训技巧在新 30 局上为 43.33，未超过单纯续训对照的 53.00。长回报原奖励模型在另一组 50 局上为 88.54，辅助奖励为 57.34；Balanced 在后续 50 局上为 100.76、标准更新对照为 81.84，但配对均分差的区间跨零。各轮测试开局不同，只记录实验过程，不直接计算历轮提升率。

双学习率实验提高价值学习率并增加额外拟合。在已记录的更新中，96.9\% 的额外拟合降低训练 MSE，相对降幅中位数约 19\%，但按开发均分选中的仍为起点。独立测试均分 128.12 与起点相同，Balanced 对照为 134.00。更低训练误差并未自动转化为更好策略。

该阶段稳定候选的开发均分为 158.53，起点为 158.59，两者十分接近；中位数却由 80.5 升至 132.5。这促使项目在新验证集上复核候选并增加均分约束的稳定性选择，而不是按最终测试重新挑模型。随后由开发流程选中的 Balanced 检查点成为 GAE 两组的共同起点。

GAE 阶段主候选由开发均分选定为 $\lambda=0.98$；$\lambda=0.99$ 作为补充消融保留。统一测试同时纳入两组的共同起点，具体分数、配对区间和截断情况在第\ref{sec:comparison}节集中分析。

\subsubsection{多次配对续训与选模稳定性}
图\ref{fig:ppo-multiseed}逐次展示 KL 对照的开发均分和中位数，保留全部实际评估节点，不把三次训练合成一条平滑的提升曲线。各次重复共享模型起点、游戏规则和开发开局，但训练采样的随机性不同；同一次重复的反馈组和对照使用相同训练随机性设置。

\reportfigure{images/ppo_multiseed_curves.png}{PPO 多次配对续训}{三次重复的开发均分、中位数和最终选中检查点。}{KL 反馈与固定计划的三次配对续训。上排为每次 32 局开发均分，下排为中位数；星号表示按开发规则最终选中的检查点，空心菱形表示均分最佳节点，二者可能重合。每列独立展示一次训练重复，未用跨重复置信带代替个体曲线。}{fig:ppo-multiseed}

曲线显示不同重复具有不同的优势阶段。重复一的反馈组在中途回落后恢复，重复二的反馈组开发均分在后期较高，重复三则在较早节点达到自身最佳，而固定计划对照在后期继续改善。由此不能把某一次、某个节点的领先解释为稳定的算法优势，也不能假定更多训练必然提高得分。

选模还需要同时考虑均分和典型表现。重复一反馈组的开发均分最高为 234.53，但该节点中位数为 185.50；依据预先约定的约束，最终选中开发均分 231.56、中位数 203.00 的稳定候选。其余重复同样依据开发数据选择，不因独立测试结果更换检查点。三次重复是同一冻结起点下的续训实验，尚不能代表从随机初始化开始的三次独立训练。

\subsubsection{改进措施的配对消融}
图\ref{fig:ppo-ablations}用“改动减同轮对照”的配对均分差汇总历轮观察。每一行的两组策略面对相同开局，零线表示样本均分相同；区间用于展示差异的不确定性。由于各行的模型起点、测试集合与训练历史不同，图中数值不构成跨轮算法排名。稳定续训和独立裁剪/Huber 均含组合改动，不能据此单独归因某一个部件；双学习率组还因早停具有不同实际预算，且最终保留起点模型，尤其不能视为纯粹的价值拟合单因素收益。

\reportfigure{images/ppo_ablation_effects.png}{PPO 改进措施消融}{各轮相对自身对照的配对收益及不确定性。}{各轮改动相对同轮对照的配对均分差。前五行展示固定策略的回合配对区间，最后一行为 KL 三次续训的训练与游戏两维配对区间；灰点表示三次 KL 续训各自的均分差。不同区间的统计对象不同，不能把跨零与否用作统一的算法优劣检验。}{fig:ppo-ablations}

在各轮固定策略样本中，辅助奖励相对原始奖励的均分差为 \ShapingAblationDiff，区间为 $[\ShapingAblationLow,\ShapingAblationHigh]$，该轮未支持其更优。稳定续训组合、独立裁剪/Huber、双学习率组合和 GAE 调整的区间均跨零，当前证据不足以认定它们具有稳定的额外收益。保留这些结果比只展示有利的中间节点更能反映实际优化过程。

KL 阶段三次续训的测试均分分别取平均，反馈组为 \KLAdaptiveMean，对照组为 \KLControlMean；逐次配对差分别为 $+\KLRepeatOneDiff$、$\KLRepeatTwoDiff$ 和 $+\KLRepeatThreeDiff$。总体配对差为 \KLAblationDiff，按训练重复对和共同游戏开局两个维度进行配对重采样，得到 95\% 区间 $[\KLAblationLow,\KLAblationHigh]$。三次重复中既有提高也有退步，区间跨零，尚不能确认 KL 反馈具有稳定优势。

这里每种设置虽然共测试 150 局，但只有三次续训重复，不能把 150 局当成 150 次独立训练。前几轮区间仅描述固定策略在回合样本上的条件差异；KL 区间额外考虑当前三个续训重复，仍受重复数量少及共享起点的限制。原始历史区间按已有统计保留，新补充的辅助奖励配对区间与 KL 两维区间均使用 20000 次重采样。

\subsubsection{KL 控制过程与更新诊断}
图\ref{fig:ppo-kl-diagnostics}将倍率、实际策略学习率、整轮 KL、平滑 KL、概率比越界比例及裁剪范围按日志记录位置对齐。三次重复全部展示，颜色区分两种设置，线型区分重复。优化指标的日志输出位置与开发评估节点并非同一采样时点，因此这些曲线用于诊断更新过程，不作逐点因果归因。

\reportfigure{images/ppo_kl_diagnostics.png}{KL 控制诊断}{学习率、完整更新后 KL、反馈阈值与概率比裁剪调度。}{KL 反馈的实际运行日志。平滑 KL 子图的阴影表示反馈阈值带；倍率和平滑值子图展示反馈组，倍率为 1 的水平线表示固定计划对照。概率比越界比例定义为 $|\rho-1|>\epsilon$ 的样本比例，不等同于实际被裁剪的梯度贡献。}{fig:ppo-kl-diagnostics}

\FloatBarrier
控制器确实改变了更新幅度：三次反馈组在前期接近固定计划，随后逐步提高倍率，并在后半段达到 2 倍上限；因此阶段末实际策略学习率约为 $10^{-5}$，对照约为 $5\times10^{-6}$，价值分支仍沿基础计划退火。最后约 10 万步中，反馈组整轮采样 KL 的中位数约为 $1.6\times10^{-4}$，对照约为 $8.5\times10^{-5}$；概率比越界比例分别约为 6\% 和 4\%。这些观察支持反馈已生效，但不能把更大的 KL 或越界比例直接解释为更好的控制。

还可以观察到，基础学习率和裁剪范围持续下降，而反馈目标保持固定。后期平滑 KL 低于阈值带时，倍率已经达到上限，控制器无法继续补偿退火。这提示下一轮可分别研究反馈目标的阶段调度或更保守的倍率上限；当前曲线尚不足以认定调度差异就是个别重复退步的原因。结合开发曲线与独立测试的波动，应继续将参数变化、候选选择和最终得分分开评价。

\FloatBarrier
